% Drop-in replacement section for the latest fixed-packet paper.
% Insert immediately after "Exact least-prime channels" and before the
% existing prime-3 Fourier discussion.

\subsection{Exact dyadic sign-flip cells}
\label{subsec:exact-dyadic-cells}

The prime-$2$ mechanism is stronger than a channel-level heuristic.  It
compresses the raw M\"obius sequence into an exact, permanent four-slot cell
before any averaging, projection, or asymptotic estimate is applied.

\begin{theorem}[Exact dyadic sign reversal]
\label{thm:exact-dyadic-sign-reversal}
For every odd integer $a\ge1$,
\begin{equation}
 \boxed{\mu(2a)=-\mu(a).}
 \label{eq:mobius-doubling-odd}
\end{equation}
The identity includes nonsquarefree $a$: in that case both sides vanish.
\end{theorem}

\begin{proof}
If $a$ is squarefree, then $\gcd(2,a)=1$ because $a$ is odd.  Multiplicativity
of the M\"obius function and $\mu(2)=-1$ give
\[
 \mu(2a)=\mu(2)\mu(a)=-\mu(a).
\]
If $a$ is not squarefree, some prime square divides $a$, hence the same prime
square divides $2a$.  Therefore $\mu(a)=\mu(2a)=0$.
\end{proof}

Taking $a=2k+1$ gives the exact middle-slot identity
\begin{equation}
 \boxed{
 \mu(4k+2)
 =\mu\bigl(2(2k+1)\bigr)
 =-\mu(2k+1).
 }
 \label{eq:four-k-plus-two-flip}
\end{equation}
Also,
\begin{equation}
 \boxed{\mu(4k+4)=0,}
 \label{eq:fourth-slot-zero}
\end{equation}
because $4\mid4k+4$.

\begin{definition}[Complete four-slot cell]
\label{def:complete-four-slot-cell}
For $k\ge0$, define
\begin{equation}
 \mathbf C_k
 :=\bigl(
    \mu(4k+1),
    \mu(4k+2),
    \mu(4k+3),
    \mu(4k+4)
   \bigr).
 \label{eq:four-slot-cell-raw}
\end{equation}
Its scalar compression is
\begin{equation}
 C_k:=\sum_{j=1}^{4}\mu(4k+j).
 \label{eq:compressed-cell-definition}
\end{equation}
\end{definition}

\begin{theorem}[Exact cell compression]
\label{thm:exact-cell-compression}
Every complete four-slot cell has the exact form
\begin{equation}
 \boxed{
 \mathbf C_k
 =\bigl(
   \mu(4k+1),
   -\mu(2k+1),
   \mu(4k+3),
   0
  \bigr).
 }
 \label{eq:four-slot-cell-exact}
\end{equation}
Consequently,
\begin{equation}
 \boxed{
 C_k
 =\mu(4k+1)-\mu(2k+1)+\mu(4k+3).
 }
 \label{eq:compressed-cell-exact}
\end{equation}
\end{theorem}

\begin{proof}
Substitute \cref{eq:four-k-plus-two-flip,eq:fourth-slot-zero} into
\cref{eq:four-slot-cell-raw,eq:compressed-cell-definition}.
\end{proof}

The sign pattern $(+,-,+,0)$ is therefore not an expected local pattern and
not a consequence of averaging.  It is an exact algebraic identity in every
complete cell.  In particular,
\begin{equation}
 \boxed{
 p=2\text{ creates the permanent }(+,-,+,0)\text{ cell}.
 }
 \label{eq:p2-creates-cell}
\end{equation}
The coherent common component of an odd source and its double has already
cancelled before the packet norm is formed:
\begin{equation}
 \mu(a)\phi_a+\mu(2a)\phi_{2a}
 =\mu(a)\bigl(\phi_a-\phi_{2a}\bigr).
 \label{eq:doubling-column-difference}
\end{equation}
Thus the analytic problem concerns the difference of two fixed incidence
columns, not two unrestricted full columns.

\subsubsection{Prime $3$ acts on the already-compressed cell}

The three active entries in the $k$th cell are $4k+1$, $4k+2$, and $4k+3$.
Since $4\equiv1\pmod3$,
\begin{equation}
 4k+j\equiv k+j\pmod3.
 \label{eq:p3-cell-residues}
\end{equation}
The three active slots therefore occupy all three residue classes modulo $3$.

\begin{theorem}[Universal prime-$3$ activation]
\label{thm:universal-prime-three-activation}
Every complete four-slot cell contains exactly one active entry divisible by
$3$.  The affected slot cycles deterministically:
\begin{equation}
\begin{array}{c|c}
 k\bmod3 & \text{active multiple of }3\\
 \hline
 0 & 4k+3\\
 1 & 4k+2\\
 2 & 4k+1
\end{array}
\label{eq:p3-slot-cycle}
\end{equation}
\end{theorem}

\begin{proof}
The residues of the three active entries are $k+1,k+2,k$ modulo $3$, which
are pairwise distinct and exhaust $\mathbb Z/3\mathbb Z$.  Reading off the
zero residue for each value of $k\bmod3$ gives \cref{eq:p3-slot-cycle}.
\end{proof}

Accordingly,
\begin{equation}
 \boxed{
 p=3\text{ necessarily acts on exactly one of the three active channels}.
 }
 \label{eq:p3-disrupts-cell}
\end{equation}
The correct order of operations is therefore
\[
 \text{raw M\"obius values}
 \xrightarrow{\ p=2\ }
 \text{exact }(+,-,+,0)\text{ cells}
 \xrightarrow{\ p=3\ }
 \text{one compulsory channel disruption per cell}.
\]
The prime-$3$ Fourier factor discussed below acts on this exact compressed
cell; it is not an independent heuristic mask placed directly on the raw
integer sequence.

\subsubsection{Exact packet compression and boundary cost}

For an integer interval $(A,B]$, let
\begin{equation}
 \mathcal K(A,B)
 :=\left\{k\ge0:
       A<4k+1\ \text{ and }\ 4k+4\le B
   \right\}
 \label{eq:complete-cell-index-set}
\end{equation}
be the complete cells contained in the interval.  The interval is the disjoint
union of these complete cells and at most two terminal fragments.

\begin{proposition}[Complete-cell packet identity]
\label{prop:complete-cell-packet-identity}
For every finite cell-index interval $I$,
\begin{equation}
 \boxed{
 \sum_{n\in4I+\{1,2,3,4\}}\mu(n)
 =\sum_{k\in I}C_k.
 }
 \label{eq:complete-cell-interval-identity}
\end{equation}
More generally,
\begin{equation}
 \sum_{A<n\le B}\mu(n)
 =\sum_{k\in\mathcal K(A,B)}C_k+E(A,B),
 \label{eq:arbitrary-packet-compression}
\end{equation}
where $E(A,B)$ is supported only on the two terminal fragments.  For arbitrary
integer endpoints, $E(A,B)$ contains at most six M\"obius values and therefore
\begin{equation}
 |E(A,B)|\le6.
 \label{eq:universal-boundary-six}
\end{equation}
\end{proposition}

For the half-square packet grid used in this paper,
\[
 X_m=\left\lfloor\frac{m^2}{2}\right\rfloor
 \equiv0\text{ or }2\pmod4.
\]
If an endpoint is $2$ modulo $4$, one of the two adjacent boundary integers is
a multiple of $4$ and has M\"obius value zero.  Hence the sharper packet bound
is available.

\begin{theorem}[Half-square packet compression]
\label{thm:half-square-packet-compression}
For
\[
 \mathcal B_{u,M}=(X_u,X_{u+M}],
 \qquad
 S_{u,M}=\sum_{X_u<n\le X_{u+M}}\mu(n),
\]
one has
\begin{equation}
 \boxed{
 S_{u,M}
 =
 \sum_{k=\lceil X_u/4\rceil}^{\lfloor X_{u+M}/4\rfloor-1}
 \left[
  \mu(4k+1)-\mu(2k+1)+\mu(4k+3)
 \right]
 +E_{u,M},
 }
 \label{eq:half-square-compressed-packet}
\end{equation}
where $E_{u,M}$ contains at most three nonzero M\"obius values.  In particular,
\begin{equation}
 \boxed{|E_{u,M}|\le3.}
 \label{eq:half-square-boundary-three}
\end{equation}
Consequently,
\begin{equation}
 \boxed{
 \sum_{u=M}^{2M-1}|E_{u,M}|^2\le9M.
 }
 \label{eq:boundary-energy-nine-M}
\end{equation}
\end{theorem}

\begin{proof}
The complete cells contained in $(X_u,X_{u+M}]$ give the displayed compressed
sum by \cref{thm:exact-cell-compression}.  Each endpoint is $0$ or $2$ modulo
$4$.  Endpoint class $0$ produces no fragment.  Endpoint class $2$ produces a
fragment of two integers, but at the left endpoint one is the terminal
multiple of $4$ and hence has M\"obius value zero.  Across both ends there are
at most three nonzero terms.  The energy estimate follows by summing the
pointwise bound $|E_{u,M}|^2\le9$ over exactly $M$ starting indices.
\end{proof}

Define the exact compressed packet
\begin{equation}
 \widetilde S_{u,M}
 :=
 \sum_{k=\lceil X_u/4\rceil}^{\lfloor X_{u+M}/4\rfloor-1}
 \left[
  \mu(4k+1)-\mu(2k+1)+\mu(4k+3)
 \right].
 \label{eq:compressed-packet-definition}
\end{equation}
Then $S_{u,M}=\widetilde S_{u,M}+E_{u,M}$ with boundary energy $O(M)$.
The raw and compressed variance criteria are therefore equivalent at the
$M^{3+\varepsilon}$ scale.  Indeed, applying
$|x+y|^2\le2|x|^2+2|y|^2$ in both directions gives
\begin{equation}
 \sum_{u=M}^{2M-1}|S_{u,M}|^2
 \asymp_{\mathrm{target}}
 \sum_{u=M}^{2M-1}|\widetilde S_{u,M}|^2,
 \label{eq:raw-compressed-energy-equivalence}
\end{equation}
where the discrepancy is bounded using \cref{eq:boundary-energy-nine-M}.

The principal packet problem should therefore be stated in its algebraically
compressed form:
\begin{equation}
 \boxed{
 \sum_{u=M}^{2M-1}
 \left|
 \sum_{k=\lceil X_u/4\rceil}^{\lfloor X_{u+M}/4\rfloor-1}
 \left[
  \mu(4k+1)-\mu(2k+1)+\mu(4k+3)
 \right]
 \right|^2
 \ll_{\varepsilon}M^{3+\varepsilon}.
 }
 \label{eq:compressed-master-variance}
\end{equation}
Thus the open frame estimate acts on an already compressed signed process.
It does not begin with raw independent M\"obius columns.

\begin{conjecture}[Compressed actual-start signed-frame contraction]
\label{conj:compressed-actual-start-frame}
For every $\varepsilon>0$ there exist $\eta>0$, $0<\rho<1$, and
$C_{\varepsilon}<\infty$, independent of $M$, such that the fixed actual-start
compressed packet state satisfies
\begin{equation}
 \mathcal E_{\mathrm{comp}}(M)
 \le
 \rho\,\mathcal E_{\mathrm{parent,comp}}(M)
 +C_{\varepsilon}M^{3-\eta+\varepsilon},
 \label{eq:compressed-frame-recursion}
\end{equation}
where the parent energy is assembled only from strictly smaller numerical
scales and all actual starts, cross-cell terms, Fejer weights, and completion
terms are retained in the same signed quadratic form.
\end{conjecture}

By the uniform closure theorem, \cref{conj:compressed-actual-start-frame}
would imply a finite $M$-independent bound after cubic normalization.  The
structural dependency is now explicit:
\begin{equation}
 \boxed{
 \mu(2a)=-\mu(a)
 \Longrightarrow
 (+,-,+,0)\text{ exact cells}
 \Longrightarrow
 p=3\text{ universal disruption}
 \Longrightarrow
 2ab\text{ displacement/lifetime control}
 \Longrightarrow
 \text{compressed signed-frame problem}.
 }
 \label{eq:structural-dependency-chain}
\end{equation}
